\icode{N24X-MG-MU-L3-261021-002}
\nceeID[5]{2024湖南}{7}
（多选）
2024年5月3日，“嫦娥六号”探测器顺利进入地月转移轨道，正式开启月球之旅。相较于“嫦娥四号”和“嫦娥五号”，本次的主要任务是登陆月球背面进行月壤采集，并通过升空器将月壤转移至绕月运行的返回舱，返回舱再通过返回轨道返回地球。设返回舱绕月运行的轨道为圆轨道，半径近似为月球半径。已知月球表面重力加速度约为地球表面的 \(\frac{1}{6}\)，月球半径约为地球半径的 \(\frac{1}{4}\)。关于返回舱在该绕月轨道上的运动，下列说法正确的是\paren[]
\begin{choices}
	\item 其相对于月球的速度大于地球第一宇宙速度
	\item 其相对于月球的速度小于地球第一宇宙速度
	\item 其绕月飞行周期约为地球上近地圆轨道卫星周期的\(\frac{2}{3}\)倍
	\item 其绕月飞行周期约为地球上近地圆轨道卫星周期的\(\frac{3}{2}\)倍
\end{choices}
\begin{answer}
	BD\qquad 提示：第一宇宙速度等于近地轨道的环绕速度
\end{answer}